\chapter{Introduzione}
\label{chap:introduzione}

\section{Contesto e motivazione}
\label{sec:intro-contesto}

La disponibilità di immagini astronomiche rende possibile studiare le
popolazioni galattiche su campioni estesi, ma richiede strumenti capaci
di organizzare e interpretare grandi quantità di osservazioni. La
morfologia costituisce uno degli elementi di questa analisi: distinguere
strutture regolari, dischi e sistemi perturbati contribuisce alla
caratterizzazione delle galassie e allo studio della loro evoluzione.
L'ispezione visuale resta un riferimento importante, ma la crescita dei
cataloghi motiva lo sviluppo di procedure automatiche.

Galaxy Zoo ha mostrato come la partecipazione di numerosi classificatori
possa produrre annotazioni morfologiche utilizzabili per analisi su larga
scala \cite{lintott2011galaxyzoo}. Su queste basi, le reti neurali
convoluzionali consentono di apprendere direttamente dalle immagini
rappresentazioni utili al riconoscimento \cite{dieleman2015galaxy}.
Il problema non si esaurisce però nel raggiungimento di un punteggio
elevato: occorre comprenderne la stabilità, la dipendenza dai dati
disponibili e il costo necessario per ottenerlo.

In questo contesto si colloca l'interesse per il quantum machine
learning. I modelli ibridi combinano elaborazioni classiche e circuiti
quantistici parametrizzati. Il transfer learning offre un collegamento
pratico tra immagini ad alta dimensionalità e registri quantistici
di piccole dimensioni: una rete preaddestrata estrae le feature e una
testa ibrida le utilizza per classificare \cite{mari2020transfer}.
La possibilità di costruire un simile modello non dimostra tuttavia
che il circuito aggiunga un beneficio rispetto a un'alternativa classica.

In questa tesi verifico il contributo della componente quantistica
alla classificazione morfologica delle galassie. Il confronto con
teste classiche permette di valutarne i possibili benefici insieme
ai costi e ai limiti, senza assumere in partenza un vantaggio.

\section{Problema affrontato e obiettivo del lavoro}
\label{sec:intro-obiettivo}

Il dataset finale comprende \(5\,628\) immagini, distribuite equamente
tra galassie ellittiche, spirali e sistemi etichettati come merger.
Si definiscono un task binario, limitato a ellittiche e spirali, e un
task a tre classi che include anche i merger. Questa seconda
configurazione permette di esaminare il comportamento dei modelli in
presenza di una categoria associata a morfologie perturbate e potenzialmente
più ambigue. Non si assume che il passaggio tra task isoli un effetto
causale della sola classe aggiunta.

Tutte le architetture utilizzano le 512 feature prodotte da una ResNet18
preaddestrata e congelata. Il confronto riguarda quindi le teste di
classificazione, non backbone addestrati separatamente. Accanto a una
testa lineare e a un percettrone multistrato, si considera un controllo
classico con bottleneck di quattro o otto componenti, denominato B2.
Il modello Q inserisce un circuito quantistico addestrabile dopo una
proiezione della stessa dimensione; QF mantiene fissi i parametri
del circuito lasciando addestrabili gli strati classici circostanti.

La domanda centrale è se l'inserimento del circuito migliori la
classificazione rispetto al controllo classico e se il suo eventuale
contributo dipenda dalla difficoltà del task e dalla disponibilità
dei dati. La scelta di baseline e ablazioni è parte essenziale
della domanda scientifica, non un semplice supporto alla graduatoria
finale \cite{bowles2024benchmarking}.

\section{Domande di ricerca}
\label{sec:intro-rq}

Il lavoro è organizzato attorno a cinque domande:
\begin{description}
 \item[RQ1 --- Confronto principale.] Il modello ibrido Q migliora
 la Macro-F1 rispetto a B2 con bottleneck della stessa dimensione?
 \item[RQ2 --- Complessità del task.] Come cambia il divario Q--B2
 passando dalla classificazione binaria a quella a tre classi,
 e quale ruolo descrittivo hanno gli errori sulla classe merger?
 \item[RQ3 --- Disponibilità dei dati.] Come variano le prestazioni
 e il confronto Q--B2 utilizzando 100, 300 e 900 immagini di training
 rispetto al training set completo?
 \item[RQ4 --- Contributo del circuito.] Quale beneficio si osserva
 addestrando i parametri quantistici rispetto a mantenerli congelati?
 Come cambiano le prestazioni passando da quattro a otto dimensioni/qubit?
 \item[RQ5 --- Costo computazionale.] Quale compromesso tra prestazioni
 e tempi di training e inferenza emerge rispetto al controllo classico?
\end{description}

Queste domande non hanno tutte lo stesso ruolo inferenziale. Il confronto
Q--B2 sul Full costituisce la famiglia primaria, quello sui subset la
famiglia secondaria. Le analisi per classe, Q--QF, dimensionalità,
calibrazione e costo completano il quadro con funzione esplorativa,
senza moltiplicare indiscriminatamente i test tra tutte le architetture.

\section{Approccio sperimentale e contributi}
\label{sec:intro-contributi}

La valutazione utilizza dieci ripetizioni di holdout stratificato.
Entro ogni ripetizione, i subset ridotti sono estratti dal training
set e condividono validation e test. La matrice comprende due task,
otto configurazioni, quattro regimi di dati e dieci seed, per un
totale di \(640\) run. Il termine Full indica l'intero training
set della ripetizione, non l'impiego di tutte le immagini del dataset
per addestrare il modello.

La Macro-F1 è la metrica principale, affiancata da metriche per classe,
qualità delle probabilità e tempi. Il confronto appaiato permette di
valutare differenze entro la stessa ripetizione. Il bootstrap stima
l'incertezza degli effetti e la correzione di Holm gestisce i confronti multipli,
con le limitazioni dovute al riuso del dataset esplicitate nel metodo.

Il lavoro comprende la costruzione del dataset, l'implementazione
dei modelli e la loro valutazione. La pipeline conserva il collegamento
tra immagini, partizioni e predizioni, così da rendere verificabili
i risultati. Il confronto tra teste permette di studiare la
compressione, l'inserimento del circuito e l'addestramento dei suoi
parametri attraverso controlli dedicati. L'analisi su due task e quattro
regimi di dati permette
poi di osservare come cambino gli errori e il costo della simulazione.

Le esecuzioni quantistiche utilizzano un simulatore classico, non un
processore quantistico fisico. Il lavoro valuta pertanto le architetture
nel protocollo considerato e non misura uno speedup su QPU. Analogamente,
un risultato su questo campione non viene interpretato come una
conclusione universale sul quantum machine learning.

\section{Organizzazione della tesi}

Il \cref{chap:fondamenti} introduce il contesto astronomico, il transfer
learning e i fondamenti delle reti ibride. Il \cref{chap:dataset}
descrive la selezione dei dati, le etichette e il protocollo di
partizionamento. Il \cref{chap:modelli} specifica architetture,
addestramento e strumenti di valutazione. Il \cref{chap:risultati}
presenta le evidenze quantitative; il \cref{chap:discussione} ne
interpreta il significato e le limitazioni. Infine, il
\cref{chap:conclusioni} riassume i contributi e delinea gli sviluppi
futuri.
